% Independently audited by random_hosts and path_stability.
% Exact checks: four-regular-terminal-avoidance-checks.json.
\section{Avoiding one terminal of a saturated pair}
\label{sec:saturated-terminal-avoidance}

\begin{lemma}[At least half the paths avoid the other terminal]
\label{lem:saturated-terminal-avoidance}
Let $T$ be a normalized three-pole of rank $h$, with distinct ports $a,b,c$.
Suppose $P_{ac}(T)=2^h$.
Then
\[
 P_{bc}(T-a)\ge\tfrac12P_{bc}(T),\qquad
 P_{ab}(T-c)\ge\tfrac12P_{ab}(T).
\]
\end{lemma}

\begin{proof}
It suffices to prove the first inequality.
Suppress $b$ and use the saturated two-pole classification.
Its positive-rank bridgeless pieces occur in series from $a$ to $c$.
If $b$ subdivides a bridge, every $b$--$c$ path avoids $a$.
The same conclusion holds when its active piece is not the first piece.
Indeed, entering an earlier piece and returning would repeat its separating bridge.

Suppose the active piece is the first piece.
Later pieces contribute the same multiplicative factor to both path counts.
We therefore work inside the active piece.
For a prefix with $p$ ladder positions, let $A_p$ be its cap count.
Let $T_p$ count prefix caps that pass through $a$.
Equivalently, $T_p$ counts the two through paths before attaching the initial cap.
The empty prefix has $A_0=T_0=1$.
A singleton block replaces $(A,T)$ by $(A+1,T)$.
A square block replaces it by $(2A+2,2T)$.
Consequently,
\begin{equation}\label{eq:prefix-cap-through-dominance}
 A_p\ge2T_p-1.
\end{equation}

Use the same prefix and suffix parameters as in the local subdivision table.
For an interface edge, a singleton rung, and a square edge, respectively, the total local $b$--$c$ counts are
\[
 (1+A_p)2^q,\qquad 2(1+A_p)2^q,\qquad (4+3A_p)2^q.
\]
The corresponding counts through $a$ are
\[
 T_p2^q,\qquad 2T_p2^q,\qquad 3T_p2^q.
\]
These follow from the same local path choices in the subdivision table.
Only the terms that use a prefix cap can meet $a$.
Replacing its cap count $A_p$ by $T_p$ gives precisely those terms.
Equation~\eqref{eq:prefix-cap-through-dominance} bounds each through-$a$ count by half its total count.

If the active piece is a digon, subdividing either edge gives two local paths.
Exactly one passes through $a$.
This also gives the required factor one half.
The reflected argument proves the second inequality.
\end{proof}

\begin{proposition}[Forest concentration beside a full $K_{2,2}$ profile]
\label{prop:five-pole-k22-forest-concentration}
Let $X$ be a normalized five-pole of rank $h$.
Write its ports as $L=\{l_1,l_2\}$, $R=\{r_1,r_2\}$, and $b$.
Suppose all four $L$--$R$ path counts equal $2^h$.
For $j\in L$, let $D_j$ count forests on odd set $\{b,j,r_1,r_2\}$ whose pairing differs from $bj\mid r_1r_2$.
For $j\in R$, define $D_j$ with $L$ and $R$ exchanged.
Then
\[
 \sum_{j\in L}D_j\le\sum_{j\in L}P_{bj}(X),\qquad
 \sum_{j\in R}D_j\le\sum_{j\in R}P_{bj}(X).
\]
\end{proposition}

\begin{proof}
Add vertices $z,w$ adjacent to the two ports of $L,R$, respectively.
This gives a normalized three-pole $T$ of rank $h+2$, with ports $z,w,b$.
Its $z$--$w$ path count is the sum of the four full cross counts.
Thus $P_{zw}(T)=2^{h+2}$.

A $b$--$z$ path that avoids $w$ corresponds to a $b$--$j$ path in $X$, for exactly one $j\in L$.
A $b$--$z$ path through $w$ corresponds to one of the forests counted by $D_j$.
Deleting the two cap vertices establishes both correspondences bijectively.
Therefore
\[
 P_{bz}(T-w)=\sum_{j\in L}P_{bj}(X),\qquad
 P_{bz}(T)=\sum_{j\in L}\bigl(P_{bj}(X)+D_j\bigr).
\]
Lemma~\ref{lem:saturated-terminal-avoidance} gives the first asserted inequality.
Exchange $L$ and $R$ to obtain the second.
\end{proof}

\begin{corollary}[An isolated fifth port cannot give an aligned critical profile]
\label{cor:five-pole-k22-isolate-deficit}
Consider two rank-$h,k$ five-poles with the same four full $K_{2,2}$ path pairs.
Suppose every path count involving the fifth port is at most $\delta$ times its shore's affine-space size.
For fixed $\delta>0$, their normalized five-cut crossing contribution is at most
\[
 4+8\delta+4\delta^2+o(1)
\]
as $\min(h,k)\longrightarrow\infty$.
In particular, this is bounded below the one-half threshold $8$ when $\delta$ is sufficiently small.
\end{corollary}

\begin{proof}
The four full path products contribute $4$.
The two same-side path densities tend to zero by the two-high-pairs tripole theorem.
The four fifth-port products contribute at most $4\delta^2$.

The forest sector on $L\cup R$ tends to zero by the high-cross-sum capping bound.
In each of the other four sectors, let $D,E$ be the normalized masses outside the distinguished cap pairing on the two shores.
The exact gluing expression for this sector is at most $D+E$.
Indeed, its subtracted equal-pairing term includes the product of the two distinguished masses.
All total forest masses are at most one.
Proposition~\ref{prop:five-pole-k22-forest-concentration} bounds the sum of these errors by $4\delta$ on each shore.
Adding the path and forest contributions proves the result.
\end{proof}

The concentration statement does not require the four forest sectors themselves to be small.
It identifies the common pairings whose products cancel in the exact five-cut identity.
